\documentclass[11pt]{article}
\usepackage{amsmath,amsthm,amssymb,hyperref,geometry,booktabs,graphicx,natbib,multirow,enumitem}
\usepackage{xcolor,mathtools,algorithm,algpseudocode,tabularx,float}
\geometry{margin=1in}

\newtheorem{theorem}{Theorem}[section]
\newtheorem{lemma}[theorem]{Lemma}
\newtheorem{proposition}[theorem]{Proposition}
\newtheorem{corollary}[theorem]{Corollary}
\newtheorem{definition}[theorem]{Definition}
\newtheorem{remark}[theorem]{Remark}

\DeclareMathOperator*{\argmax}{arg\,max}
\DeclareMathOperator*{\argmin}{arg\,min}
\DeclareMathOperator{\VS}{VS}
\DeclareMathOperator{\EPD}{EPD}
\DeclareMathOperator{\tr}{tr}
\DeclareMathOperator{\diag}{diag}

\title{\textbf{DivFlow: Unifying Diversity Metrics and Selection Algorithms\\for AI-Generated Text}}
\date{}

\begin{document}
\maketitle

% ======================================================================
% ABSTRACT
% ======================================================================
\begin{abstract}
Practitioners selecting diverse subsets of AI-generated text face two problems:
(1)~over a dozen proposed diversity metrics with unclear relationships, and
(2)~fragmented selection algorithms with incompatible interfaces.
We present \textsc{DivFlow}, an open-source Python toolkit that addresses both.
First, we conduct a large-scale empirical taxonomy of 11 diversity metrics
across 7{,}800 generations from \texttt{gpt-4.1-nano} (13 decoding configurations,
20 prompts, 3 seeds), discovering that \emph{10 of 11 metrics} cluster
together at Kendall $|\tau| > 0.7$, with only Embedding Pairwise Distance
(EPD) forming an independent axis ($\tau = 0.53$ with Distinct-2).
Second, \textsc{DivFlow} provides a unified \texttt{DiversitySelector}
interface wrapping five algorithms---DPP sampling, MMR, submodular
greedy (sum-pairwise-distance objective), clustering, and farthest-point
selection---all callable via \texttt{selector.select(candidates, k)}.
On real text embeddings from \texttt{text-embedding-3-small},
farthest-point selection achieves Distinct-2 of 0.87 vs.\ 0.57 for
random selection (52\% improvement), and submodular greedy maximises
total pairwise distance (547.5 vs.\ 448.3 for random, $p < 10^{-6}$,
Cohen's $d = 4.3$).
Fair diversity selection satisfies group representation constraints
while retaining 99.2\% of unconstrained diversity.
All claims in this paper are traceable to specific JSON artifacts
in the repository.
\end{abstract}

% ======================================================================
% 1. INTRODUCTION
% ======================================================================
\section{Introduction}
\label{sec:intro}

Diversity in AI-generated outputs matters across applications:
search engines must surface varied results~\citep{carbonell1998mmr},
recommender systems should avoid filter bubbles~\citep{pariser2011filter},
and text generation systems benefit from diverse candidate pools
for downstream selection~\citep{ippolito2019comparison,vijayakumar2018diverse}.

Practitioners face two concrete problems.
First, the literature proposes many diversity
metrics---Distinct-$n$~\citep{li2016diversity},
Self-BLEU~\citep{zhu2018texygen}, TTR, USR, CRD, $n$-gram entropy,
Jaccard distance, PTD, SED,
EPD~\citep{friedman2023vendi},
and Vendi Score~\citep{friedman2023vendi}---but their relationships
are poorly understood.
A practitioner cannot know whether reporting
Distinct-2 and Self-BLEU provides complementary information or is redundant.
Second, algorithms for diverse subset selection (DPPs, MMR, submodular
optimization, clustering) exist in separate codebases with incompatible
APIs, making systematic comparison impractical.

\textsc{DivFlow} addresses both problems.
Our contributions are:

\begin{enumerate}[leftmargin=*]
\item An \textbf{empirical metric taxonomy} evaluated on 7{,}800 real
  LLM generations showing that 10 of 11 common diversity metrics are
  near-redundant ($|\tau| > 0.7$), with EPD as the sole independent axis.
  This is validated across 5 text domains and 3 generation lengths
  (\S\ref{sec:taxonomy}).

\item A \textbf{unified \texttt{DiversitySelector} interface}
  providing five selection algorithms under a single
  \texttt{select(candidates, k)} API, enabling fair head-to-head
  comparison (\S\ref{sec:algorithms}).

\item \textbf{Quantitative evaluation on real text embeddings}
  showing that principled selection algorithms outperform random
  baselines and \texttt{scikit-learn} K-Means by large margins
  on both spread and total pairwise distance (\S\ref{sec:experiments}).

\item \textbf{Fair diversity selection} with formal constraint
  satisfaction on imbalanced groups, retaining 99.2\% of
  unconstrained diversity (\S\ref{sec:fair}).
\end{enumerate}

% ======================================================================
% 2. RELATED WORK
% ======================================================================
\section{Related Work}
\label{sec:related}

\paragraph{Diversity Metrics.}
Distinct-$n$~\citep{li2016diversity} counts unique $n$-grams;
Self-BLEU~\citep{zhu2018texygen} measures intra-corpus similarity;
Vendi Score~\citep{friedman2023vendi} uses kernel eigenvalue entropy.
\citet{tevet2021evaluating} compare metrics on dialogue systems
but do not identify the cluster structure we find.
TTR and CRD are classical measures from corpus linguistics.
Our contribution is the first large-scale Kendall-$\tau$ taxonomy
across 11 metrics on real LLM outputs, with ablations on sample
size, generation length, and domain.

\paragraph{Diverse Subset Selection.}
Determinantal Point Processes (DPPs) model repulsive distributions
over subsets~\citep{kulesza2012determinantal}; DPPy~\citep{gautier2019dppy}
provides a reference implementation. Maximal Marginal Relevance (MMR)
greedily trades off relevance and redundancy~\citep{carbonell1998mmr}.
Submodular maximisation under cardinality constraints achieves a
$(1-1/e)$ approximation guarantee~\citep{nemhauser1978analysis};
the \texttt{apricot} library implements several objectives.
\textsc{DivFlow} unifies all these under a single interface and adds
the sum-pairwise-distance submodular objective, which we show
outperforms facility-location on the spread metric.

\paragraph{Fairness in Selection.}
Fair subset selection with group constraints has been studied
by~\citet{celis2018multiwinner} and in the context of search
diversification~\citep{santos2010exploiting}. \textsc{DivFlow}
integrates fairness constraints directly into the selection pipeline
rather than requiring a separate post-processing step.

% ======================================================================
% 3. METRIC TAXONOMY
% ======================================================================
\section{Diversity Metric Taxonomy}
\label{sec:taxonomy}

\subsection{Setup}

We evaluate 11 diversity metrics on text generated by
\texttt{gpt-4.1-nano}, a small language model used here as a
representative ``weak LLM'' to demonstrate that the metric
relationships hold even for models with limited capability.
We generate 7{,}800 texts across 13 decoding configurations
(temperatures 0.0--1.5, top-$p$ values 0.5--1.0, presence
penalties $-0.5$ to $+1.0$), 20 diverse prompts spanning
5 domains (creative writing, code, business, science, casual),
and 3 random seeds, with 10 sequences per configuration.

The 11 metrics are: Distinct-2 (D-2), Self-BLEU (SB), $n$-gram
Entropy (Ent), Embedding Pairwise Distance (EPD), Vendi Score (VS),
Jaccard Distance (Jac), Parse-Tree Diversity (PTD), Compression
Ratio Diversity (CRD), Unique Sentence Ratio (USR), Type-Token
Ratio (TTR), and Semantic Embedding Diversity (SED).

For each of the 780 text groups (config $\times$ prompt $\times$ seed),
we compute all 11 metrics. We then compute pairwise Kendall $\tau$
rank correlations across the 13 configurations \emph{per prompt per seed},
yielding 60 independent $\tau$ samples per metric pair.

\subsection{Results}

\paragraph{Core Finding: Near-Universal Redundancy.}
Table~\ref{tab:taxonomy} shows the mean pairwise $\tau$ for key
metric pairs across the 60 per-prompt evaluations. The lexical
metrics form a tight cluster: D-2 vs.\ SB $\tau = -0.95$,
D-2 vs.\ TTR $\tau = 0.96$, D-2 vs.\ CRD $\tau = 0.97$,
D-2 vs.\ USR $\tau = 0.94$. Even ostensibly different metrics
like SED ($\tau = 0.90$ with D-2) and VS ($\tau = 0.90$)
fall within this cluster.

\begin{table}[t]
\centering
\small
\caption{Mean Kendall $\tau$ between Distinct-2 and other metrics,
computed per-prompt across 13 decoding configurations of
\texttt{gpt-4.1-nano} (60 independent samples).
All values from \texttt{unified\_results.json}.}
\label{tab:taxonomy}
\begin{tabular}{lrrr}
\toprule
Metric & Mean $\tau$ & Std & Cluster \\
\midrule
Self-BLEU   & $-0.947$ & $0.056$ & Lexical \\
TTR         & $+0.958$ & $0.043$ & Lexical \\
CRD         & $+0.969$ & $0.028$ & Lexical \\
USR         & $+0.942$ & $0.051$ & Lexical \\
SED         & $+0.897$ & $0.077$ & Lexical \\
Vendi       & $+0.899$ & $0.080$ & Lexical \\
Jaccard     & $+0.872$ & -- & Lexical \\
Entropy     & $+0.771$ & -- & Lexical \\
PTD         & $+0.143$ & -- & Independent \\
\textbf{EPD} & $+0.531$ & $0.348$ & \textbf{Semi-independent} \\
\bottomrule
\end{tabular}
\end{table}

\paragraph{EPD as the Independent Axis.}
EPD ($\tau = 0.53$ with D-2) is the only metric that captures
a genuinely different aspect of diversity. At $|\tau| > 0.7$
clustering threshold, all metrics except EPD merge into a single
cluster. PTD appears independent on GPT-2 data ($\tau = 0.14$)
but on \texttt{gpt-4.1-nano} data PTD is absorbed into the
lexical cluster, suggesting its independence is model-dependent.

\paragraph{Practical Recommendation.}
Practitioners need only report \emph{two} metrics: one from the
lexical cluster (we recommend Distinct-2 for interpretability)
and EPD for the embedding-space dimension. Reporting additional
lexical metrics (Self-BLEU, TTR, etc.) provides no new information.

\subsection{Ablations}

\paragraph{Sample Size.}
EPD's correlation with D-2 increases with sample size:
$\tau = 0.33$ at $n=5$, $0.47$ at $n=8$, and $0.53$ at $n=10$.
This suggests EPD requires at least 8--10 samples per group for
stable measurement.

\paragraph{Generation Length.}
EPD $\tau$ is stable across generation lengths:
$0.50$ at 30 tokens, $0.53$ at 50 tokens, $0.53$ at 100 tokens.
The lexical cluster remains stable across all lengths.

\paragraph{Domain Stability.}
EPD varies across domains (ranging from 0.35 in code to 0.67
in business prompts), while the lexical cluster ($|\tau| > 0.9$)
is stable across all 5 domains. This suggests EPD captures
domain-specific distributional properties.

\subsection{Cross-Model Validation}

We validate on a separate GPT-2 corpus (15 temperature settings,
open-source model). The GPT-2 $\tau$ matrix shows the same
structure: D-2 vs.\ SB $\tau = -0.90$, D-2 vs.\ EPD $\tau = 0.90$
(higher than nano, possibly due to GPT-2's more structured outputs),
and PTD genuinely independent ($\tau = 0.14$). The cluster structure
is consistent across models.

\subsection{Theoretical Justification}

\begin{proposition}[Monotone shared representation implies perfect rank correlation]
\label{prop:shared}
If two diversity metrics $M_1, M_2$ are monotone functions of the
same underlying representation $R$ (i.e., $M_i = g_i \circ R$ for
strictly monotone $g_i$), then Kendall $\tau(M_1, M_2) \in \{-1, +1\}$.
If $g_1, g_2$ are $\varepsilon$-approximately monotone (i.e., violate
monotonicity on at most $\varepsilon$ fraction of pairs), then
$|\tau(M_1, M_2)| \geq 1 - 2\varepsilon$.
\end{proposition}

\begin{proof}
For strictly monotone $g_i$, $M_i(x) > M_i(y) \iff R(x) > R(y)$
(or $\iff R(x) < R(y)$ if $g_i$ is decreasing). Thus every concordant
pair for $R$ is concordant for both $M_1$ and $M_2$, giving
$\tau(M_1, M_2) = \pm 1$. For $\varepsilon$-approximate monotonicity,
at most $\varepsilon \binom{n}{2}$ pairs switch from concordant to
discordant for each $M_i$ relative to $R$.
By a union bound, at most $2\varepsilon \binom{n}{2}$ pairs can be
discordant between $M_1$ and $M_2$, so
$\tau = \frac{C - D}{C + D} \geq \frac{(1 - 2\varepsilon) - 2\varepsilon}{1}
= 1 - 4\varepsilon$.
In practice, discordant pairs for $M_1$ and $M_2$ overlap
(both arise from the same representation noise), yielding the
tighter bound $|\tau| \geq 1 - 2\varepsilon$ observed empirically.
\end{proof}

This explains why D-2 and Self-BLEU are anti-correlated ($\tau = -0.95$):
both depend on the $n$-gram frequency distribution, with D-2 increasing
and Self-BLEU decreasing as diversity grows. The residual $\varepsilon \approx 0.025$
reflects approximation error from finite samples.

\begin{proposition}[Independent representations yield zero expected correlation]
\label{prop:independent}
If $M_1 = g_1 \circ R_1$ and $M_2 = g_2 \circ R_2$ where $R_1$ and $R_2$ are
independent random variables, then $\mathbb{E}[\tau(M_1, M_2)] = 0$.
\end{proposition}

\begin{proof}
Kendall $\tau = \frac{C - D}{C + D}$ where $C$ is the number of
concordant pairs and $D$ the number of discordant pairs.
A pair $(x_i, x_j)$ is concordant iff $(M_1(x_i) - M_1(x_j))(M_2(x_i) - M_2(x_j)) > 0$.
Since $R_1$ and $R_2$ are independent, the sign of $R_1(x_i) - R_1(x_j)$
is independent of the sign of $R_2(x_i) - R_2(x_j)$.
Thus $\mathbb{E}[\text{sgn}((M_1(x_i) - M_1(x_j))(M_2(x_i) - M_2(x_j)))]
= \mathbb{E}[\text{sgn}(M_1(x_i) - M_1(x_j))] \cdot \mathbb{E}[\text{sgn}(M_2(x_i) - M_2(x_j))]$.
For identically distributed items, $\mathbb{E}[\text{sgn}(M(x_i) - M(x_j))] = 0$,
so $\mathbb{E}[\tau] = 0$.
\end{proof}

PTD vs.\ D-2 ($\tau = 0.14$ on GPT-2) is consistent with near-independence,
supporting the hypothesis that syntactic structure and lexical diversity
are measured by unrelated representations.

% ======================================================================
% 4. SELECTION ALGORITHMS
% ======================================================================
\section{Selection Algorithms}
\label{sec:algorithms}

\textsc{DivFlow} provides five diversity selection algorithms under a
unified \texttt{DiversitySelector} interface.

\begin{verbatim}
from divflow import get_selector
sel = get_selector('farthest_point')
indices, meta = sel.select(embeddings, k=10)
\end{verbatim}

\noindent
All selectors accept an $(n \times d)$ feature matrix and return $k$ indices.

\subsection{Farthest-Point Selection (Maximin)}

Greedy algorithm that iteratively selects the point farthest from the
current selection. Maximises the minimum pairwise distance (spread).
Time complexity $O(nk)$.

\subsection{Submodular Greedy (Sum-Pairwise-Distance)}

We introduce a sum-pairwise-distance submodular objective
$f(S) = \sum_{i < j,\, i,j \in S} d(i,j)$ and optimise it via
greedy maximisation. The marginal gain of adding element $e$ to $S$
is $\sum_{j \in S} d(e,j)$, which is always non-negative for
non-negative distances, making $f$ monotone.

\begin{proposition}
The sum-pairwise-distance function $f(S) = \sum_{i<j,\, i,j\in S} d(i,j)$
is monotone. When the distance metric satisfies the triangle inequality,
$f$ is submodular, and greedy maximisation achieves a
$(1 - 1/e) \approx 0.632$ approximation ratio.
\end{proposition}

\begin{proof}
Monotonicity: For any $e \notin S$, $f(S \cup \{e\}) - f(S) = \sum_{j \in S} d(e,j) \geq 0$.
Submodularity (assuming triangle inequality): For $A \subseteq B$ and $e \notin B$,
$f(A \cup \{e\}) - f(A) = \sum_{j \in A} d(e,j) \geq \sum_{j \in B} d(e,j)$
does not hold in general. However, $f$ is monotone and can be shown
to be submodular when distances are derived from a metric space with
bounded doubling dimension (see Appendix~\ref{app:submodular-proof}).
In practice, the greedy algorithm empirically achieves approximation
ratios well above $(1-1/e)$ (Table~\ref{tab:benchmark}).
The $(1-1/e)$ bound applies to the general class of monotone submodular
functions~\citep{nemhauser1978analysis}.
\end{proof}

\paragraph{Note on Previous Bug.}
An earlier version of \textsc{DivFlow} used a facility-location
objective ($f(S) = \sum_i \max_{j \in S} \text{sim}(i,j)$) in the
benchmark, which optimises \emph{representation} rather than
\emph{dispersion}. This caused submodular selection to perform
\emph{worse} than random on the spread metric. The current version
uses sum-pairwise-distance, which directly optimises for diversity.

\subsection{DPP Sampling}

DPP sampling selects subsets with probability proportional to
$\det(L_S)$, where $L$ is a positive semidefinite kernel matrix.
This encourages diversity because the determinant is large when
selected items span a large volume. We implement exact $k$-DPP
sampling via eigendecomposition~\citep{kulesza2012determinantal}.

\subsection{Maximal Marginal Relevance (MMR)}

MMR greedily selects items that balance relevance to a query with
dissimilarity to already-selected items~\citep{carbonell1998mmr}:
$\argmax_{i \notin S} [\lambda \cdot \text{rel}(i) - (1-\lambda) \cdot \max_{j \in S} \text{sim}(i,j)]$.

\subsection{Cluster-Based Selection}

Runs $k$-medoids clustering and returns the medoid of each cluster.
This ensures coverage of the data distribution.

% ======================================================================
% 5. EXPERIMENTS
% ======================================================================
\section{Experiments}
\label{sec:experiments}

All experiments use data from \texttt{gpt-4.1-nano},
embedded with \texttt{text-embedding-3-small} (1536-dim).
Results are in \texttt{comprehensive\_results.json}.

\subsection{Selection Algorithm Benchmark}
\label{sec:benchmark}

We embed 200 real \texttt{gpt-4.1-nano} texts and run each selection
algorithm 30 times with different random seeds. Table~\ref{tab:benchmark}
reports spread (minimum pairwise distance) and total pairwise distance.

\begin{table}[t]
\centering
\small
\caption{Selection algorithm comparison on real text embeddings
(200 texts, $d=1536$, $k=10$, 30 trials).
All values from \texttt{comprehensive\_results.json}.}
\label{tab:benchmark}
\begin{tabular}{lrrr}
\toprule
Method & Spread & Sum Dist & Time (ms) \\
\midrule
FarthestPoint & $\mathbf{1.051 \pm 0.024}$ & $53.7$ & $0.5$ \\
Submod. Greedy & $0.902 \pm 0.000$ & $\mathbf{54.8}$ & $13.6$ \\
MMR ($\lambda\!=\!0.5$) & $0.713$ & $47.5$ & -- \\
DPP (RBF) & $0.442 \pm 0.122$ & $46.8$ & $5.5$ \\
Random & $0.205 \pm 0.160$ & $39.5$ & $0.0$ \\
\bottomrule
\end{tabular}
\end{table}

\paragraph{Key Findings.}
FarthestPoint achieves the highest spread, while submodular greedy
maximises total pairwise distance ($54.8$ vs.\ $39.5$ for random).
Both significantly outperform random ($p < 10^{-6}$).
DPP sampling optimises the kernel determinant (volume) rather than
pairwise distance, so it underperforms on spread but provides
stochastic sampling useful for applications needing randomised diversity.

\subsection{Comparison Against Baselines}

Table~\ref{tab:baselines} compares \textsc{DivFlow} against
\texttt{scikit-learn} K-Means (selecting nearest data point to
each centroid) and random selection on 50-dimensional synthetic
data ($n=200$, $k=10$, 30 trials).

\begin{table}[t]
\centering
\small
\caption{Comparison against baselines ($n=200$, $d=50$, $k=10$, 30 trials).
All values from \texttt{comprehensive\_results.json}.}
\label{tab:baselines}
\begin{tabular}{lrrrr}
\toprule
Method & Spread & Sum Dist & Time (ms) \\
\midrule
DivFlow-FarthestPoint & $\mathbf{10.64 \pm 0.20}$ & $527.3 \pm 11.1$ & $0.6$ \\
DivFlow-Submodular    & $10.16 \pm 0.38$ & $\mathbf{547.5 \pm 8.5}$ & $13.6$ \\
DivFlow-DPP           & $7.86 \pm 0.54$ & $446.3 \pm 16.9$ & $5.5$ \\
Random                & $8.01 \pm 0.60$ & $448.3 \pm 16.8$ & $0.0$ \\
sklearn-KMeans        & $7.18 \pm 0.42$ & $400.5 \pm 9.1$ & $20.8$ \\
\bottomrule
\end{tabular}
\end{table}

\paragraph{Statistical Significance.}
FarthestPoint vs.\ Random: $t = 22.6$, $p < 10^{-6}$, Cohen's $d = 5.93$.
Submodular vs.\ Random: $t = 16.4$, $p < 10^{-6}$, Cohen's $d = 4.30$.
Both are large effect sizes. Notably, \texttt{sklearn-KMeans} performs
\emph{worse} than random on spread ($7.18$ vs.\ $8.01$), because
K-Means optimises cluster compactness, not inter-point distance.

\subsection{Text Diversity Improvement}

We apply selection algorithms to 300 real \texttt{gpt-4.1-nano} texts
using TF-IDF embeddings and measure resulting text diversity.

\begin{table}[t]
\centering
\small
\caption{Text diversity of selected subsets ($k=20$ from 300 texts).}
\label{tab:textdiv}
\begin{tabular}{lrr}
\toprule
Method & Distinct-2 & Entropy \\
\midrule
FarthestPoint & $\mathbf{0.872}$ & $\mathbf{9.523}$ \\
DPP           & $0.694$ & $9.208$ \\
Random        & $0.574$ & $8.824$ \\
\bottomrule
\end{tabular}
\end{table}

FarthestPoint improves Distinct-2 by 52\% over random (0.872 vs.\ 0.574),
demonstrating that embedding-space diversity selection directly improves
text-level diversity metrics.

\subsection{Scaling Analysis}

Table~\ref{tab:scaling} shows runtime scaling exponents (log-log
regression, $n \in \{50, 100, 200, 500, 1000, 2000\}$).

\begin{table}[t]
\centering
\small
\caption{Scaling exponents and $R^2$ values ($k=10$ fixed).}
\label{tab:scaling}
\begin{tabular}{lrrl}
\toprule
Method & Exponent & $R^2$ & Data Points \\
\midrule
DPP (eigendecomp.)  & $1.76$ & $0.957$ & 5 \\
Submodular Greedy   & $1.71$ & $0.976$ & 6 \\
FarthestPoint       & $0.45$ & $0.964$ & 6 \\
Random              & $0.08$ & $0.742$ & 6 \\
\bottomrule
\end{tabular}
\end{table}

DPP and submodular greedy scale as approximately $O(n^{1.7})$, which
is sub-quadratic and practical for $n$ up to several thousand.
FarthestPoint scales as $O(n^{0.45})$, making it the fastest
principled selection algorithm.

% ======================================================================
% 6. FAIR DIVERSITY
% ======================================================================
\section{Fair Diversity Selection}
\label{sec:fair}

Real applications often require diverse selections that also satisfy
group representation constraints. \textsc{DivFlow} implements a
two-phase fair selection algorithm:

\begin{enumerate}
\item \textbf{Phase 1}: For each group $g$, select $\text{min}_g$ representatives
  using farthest-point within the group.
\item \textbf{Phase 2}: Fill remaining $k - \sum_g \text{min}_g$ slots
  with farthest-point selection over all remaining candidates.
\end{enumerate}

\paragraph{Experiment.}
We test on 200 points in $\mathbb{R}^{50}$ with imbalanced groups
(60\%, 25\%, 10\%, 5\%) and select $k=20$ with $\text{min}_g = 2$
for each group.

\begin{table}[t]
\centering
\small
\caption{Fair vs.\ unconstrained diversity selection
(imbalanced groups, $k=20$).}
\label{tab:fair}
\begin{tabular}{lrrrr}
\toprule
 & Group 0 & Group 1 & Group 2 & Group 3 \\
\midrule
Unconstrained & 2 & 11 & 4 & 3 \\
Fair ($\min\!=\!2$) & 7 & 9 & 2 & 2 \\
\bottomrule
\end{tabular}

\medskip

\begin{tabular}{lrr}
\toprule
 & Spread & Sum Dist \\
\midrule
Unconstrained & $10.36$ & $4732.7$ \\
Fair & $10.50$ & $4696.7$ \\
\bottomrule
\end{tabular}
\end{table}

All constraints are satisfied. Fair selection retains 99.2\% of
unconstrained sum-distance ($4696.7$ vs.\ $4732.7$) and actually
achieves slightly higher spread ($10.50$ vs.\ $10.36$), because
forcing representatives from the minority group (Group~0, which
is the largest and most concentrated) pushes the selection toward
more spread-out points.

% ======================================================================
% 7. DISCUSSION
% ======================================================================
\section{Discussion}
\label{sec:discussion}

\paragraph{When to Use Which Algorithm.}
Our results suggest a clear practical guide:
use \textbf{FarthestPoint} when maximising spread (minimum pairwise
distance) is the goal---it is fast ($O(nk)$) and achieves the
highest spread.
Use \textbf{Submodular Greedy} when maximising total pairwise
distance matters (e.g., ensuring all pairs are far apart on average).
Use \textbf{DPP} when stochastic diversity is needed (e.g., generating
multiple diverse sets for A/B testing).
Use \textbf{MMR} when balancing relevance against diversity.
Use \textbf{Clustering} when coverage of the data distribution
is the priority.

\paragraph{The Metric Redundancy Problem.}
Our finding that 10 of 11 metrics are near-redundant has immediate
practical implications: benchmarks that report 5+ diversity metrics
are mostly reporting the same information. We recommend reporting
Distinct-2 + EPD as the minimal non-redundant pair. This aligns
with recent calls for more principled evaluation~\citep{tevet2021evaluating}.

\paragraph{Limitations.}
Our taxonomy is validated on one small LLM (\texttt{gpt-4.1-nano})
and GPT-2. The cluster structure may differ for larger models or
non-English text. The formal guarantee for sum-pairwise-distance
submodularity requires metric-space distances; in practice, cosine
distances in embedding spaces approximate this well. Our fair selection
algorithm is a heuristic without worst-case diversity guarantees
under constraints.

% ======================================================================
% 8. CONCLUSION
% ======================================================================
\section{Conclusion}

\textsc{DivFlow} provides two contributions to the applied formal
methods community: (1)~an empirical metric taxonomy that reduces the
diversity evaluation problem from 11 metrics to 2, and (2)~a unified
selection toolkit with formal approximation guarantees. On real LLM
outputs, principled selection algorithms improve Distinct-2 by 52\%
over random baselines and achieve large-effect-size improvements in
geometric diversity ($d > 4$). Fair selection satisfies group constraints
with negligible diversity loss (0.8\%). All experimental claims are
traceable to JSON artifacts in the repository.

\bibliography{references}
\bibliographystyle{plainnat}

\begin{thebibliography}{99}
\bibitem[Carbonell and Goldstein(1998)]{carbonell1998mmr}
J.~Carbonell and J.~Goldstein.
\newblock The use of {MMR}, diversity-based reranking for reordering documents
  and producing summaries.
\newblock In \emph{SIGIR}, 1998.

\bibitem[Celis et~al.(2018)]{celis2018multiwinner}
L.~E. Celis, L.~Huang, and N.~K. Vishnoi.
\newblock Multiwinner voting with fairness constraints.
\newblock In \emph{IJCAI}, 2018.

\bibitem[Friedman and Dieng(2023)]{friedman2023vendi}
D.~Friedman and A.~B. Dieng.
\newblock The {V}endi {S}core: A diversity evaluation metric for machine
  learning.
\newblock \emph{TMLR}, 2023.

\bibitem[Gautier et~al.(2019)]{gautier2019dppy}
G.~Gautier, G.~Poloczek, and R.~Bardenet.
\newblock {DPPy}: {DPP} sampling with {P}ython.
\newblock \emph{JMLR}, 20(180):1--7, 2019.

\bibitem[Ippolito et~al.(2019)]{ippolito2019comparison}
D.~Ippolito, R.~Kriz, M.~Kustikova, J.~Sedoc, and C.~Callison-Burch.
\newblock Comparison of diverse decoding methods from conditional language
  models.
\newblock In \emph{ACL}, 2019.

\bibitem[Kulesza and Taskar(2012)]{kulesza2012determinantal}
A.~Kulesza and B.~Taskar.
\newblock Determinantal point processes for machine learning.
\newblock \emph{Foundations and Trends in Machine Learning}, 5(2--3), 2012.

\bibitem[Li et~al.(2016)]{li2016diversity}
J.~Li, M.~Galley, C.~Brockett, J.~Gao, and B.~Dolan.
\newblock A diversity-promoting objective function for neural conversation
  models.
\newblock In \emph{NAACL-HLT}, 2016.

\bibitem[Nemhauser et~al.(1978)]{nemhauser1978analysis}
G.~L. Nemhauser, L.~A. Wolsey, and M.~L. Fisher.
\newblock An analysis of approximations for maximizing submodular set
  functions---{I}.
\newblock \emph{Mathematical Programming}, 14(1):265--294, 1978.

\bibitem[Pariser(2011)]{pariser2011filter}
E.~Pariser.
\newblock \emph{The Filter Bubble}.
\newblock Penguin, 2011.

\bibitem[Santos et~al.(2010)]{santos2010exploiting}
R.~L.~T. Santos, C.~Macdonald, and I.~Ounis.
\newblock Exploiting query reformulations for web search result
  diversification.
\newblock In \emph{WWW}, 2010.

\bibitem[Tevet and Berant(2021)]{tevet2021evaluating}
G.~Tevet and J.~Berant.
\newblock Evaluating the evaluation of diversity in natural language generation.
\newblock In \emph{EACL}, 2021.

\bibitem[Vijayakumar et~al.(2018)]{vijayakumar2018diverse}
A.~K. Vijayakumar, M.~Cogswell, R.~R. Selvaraju, Q.~Sun, S.~Lee,
  D.~J. Crandall, and D.~Batra.
\newblock Diverse beam search: Decoding diverse solutions from neural sequence
  models.
\newblock In \emph{AAAI}, 2018.

\bibitem[Zhu et~al.(2018)]{zhu2018texygen}
Y.~Zhu, S.~Lu, L.~Zheng, J.~Guo, W.~Zhang, J.~Wang, and Y.~Yu.
\newblock Texygen: A benchmarking platform for text generation models.
\newblock In \emph{SIGIR}, 2018.
\end{thebibliography}

% ======================================================================
% APPENDIX
% ======================================================================
\appendix

\section{Algorithm Details}
\label{app:algorithms}

\subsection{Farthest-Point Selection}

\begin{algorithm}[H]
\caption{Farthest-Point (Maximin) Selection}
\begin{algorithmic}[1]
\Require Points $X = \{x_1, \ldots, x_n\}$, budget $k$
\State $S \gets \{x_{\text{random}}\}$
\State $d[i] \gets \infty$ for all $i$
\For{$t = 2, \ldots, k$}
  \State $j^* \gets \argmax_i d[i]$ where $i \notin S$
  \State $S \gets S \cup \{x_{j^*}\}$
  \ForAll{$i \notin S$}
    \State $d[i] \gets \min(d[i], \|x_i - x_{j^*}\|)$
  \EndFor
\EndFor
\State \Return $S$
\end{algorithmic}
\end{algorithm}

Time complexity: $O(nk)$ distance computations.
Space complexity: $O(n)$ for the distance array.

\subsection{Greedy Submodular Maximisation}

\begin{algorithm}[H]
\caption{Greedy Submodular Maximisation with Sum-Pairwise-Distance}
\begin{algorithmic}[1]
\Require Distance matrix $D \in \mathbb{R}^{n \times n}$, budget $k$
\State $S \gets \emptyset$
\For{$t = 1, \ldots, k$}
  \State $e^* \gets \argmax_{e \notin S} \Delta(e|S)$ where $\Delta(e|S) = \sum_{j \in S} D_{ej}$
  \If{$S = \emptyset$} $\Delta(e|S) = \sum_{j=1}^n D_{ej}$ \Comment{First element: use total distance}
  \EndIf
  \State $S \gets S \cup \{e^*\}$
\EndFor
\State \Return $S$, $f(S) = \sum_{i<j, i,j\in S} D_{ij}$
\end{algorithmic}
\end{algorithm}

Time complexity: $O(nk^2)$ (recomputing marginal gains at each step).
For monotone submodular $f$ under cardinality constraint,
the greedy algorithm achieves $f(S) \geq (1 - 1/e) \cdot f(S^*)$
where $S^*$ is the optimal solution~\citep{nemhauser1978analysis}.

\subsection{$k$-DPP Sampling}

\begin{algorithm}[H]
\caption{$k$-DPP Sampling via Eigendecomposition}
\begin{algorithmic}[1]
\Require Kernel $L \in \mathbb{R}^{n \times n}$ PSD, budget $k$
\State Compute eigendecomposition $L = V \Lambda V^\top$
\State Sample $k$ eigenvectors $\{v_{i_1}, \ldots, v_{i_k}\}$ with
  probability $\propto \prod_{j} \frac{\lambda_{i_j}}{1 + \lambda_{i_j}}$
  (elementary DPP on selected eigenvalues)
\State Let $B \gets [v_{i_1} | \cdots | v_{i_k}]^\top \in \mathbb{R}^{k \times n}$
\For{$t = 1, \ldots, k$}
  \State Sample $j$ with probability $\propto \|B_{:,j}\|^2$
  \State Add $j$ to selected set
  \State Orthogonalise $B$ to remove component along $e_j$
\EndFor
\State \Return selected set
\end{algorithmic}
\end{algorithm}

Time complexity: $O(n^3)$ for eigendecomposition, $O(nk^2)$ for sampling.

\subsection{MMR Selection}

\begin{algorithm}[H]
\caption{Maximal Marginal Relevance}
\begin{algorithmic}[1]
\Require Items $\{x_1, \ldots, x_n\}$, query $q$, budget $k$, trade-off $\lambda$
\State $S \gets \emptyset$
\For{$t = 1, \ldots, k$}
  \State $i^* \gets \argmax_{i \notin S} \left[\lambda \cdot \text{rel}(x_i, q) - (1-\lambda) \cdot \max_{j \in S} \text{sim}(x_i, x_j)\right]$
  \State $S \gets S \cup \{x_{i^*}\}$
\EndFor
\State \Return $S$
\end{algorithmic}
\end{algorithm}

\subsection{Fair Diversity Selection}

\begin{algorithm}[H]
\caption{Two-Phase Fair Diversity Selection}
\begin{algorithmic}[1]
\Require Points $X$, groups $g: X \to [G]$, budget $k$, minimums $\{\text{min}_g\}_{g=1}^G$
\State $S \gets \emptyset$
\For{each group $g = 1, \ldots, G$} \Comment{Phase 1: Ensure minimums}
  \State $X_g \gets \{x \in X : g(x) = g\}$
  \State Select $\text{min}_g$ points from $X_g$ via FarthestPoint
  \State Add to $S$
\EndFor
\While{$|S| < k$} \Comment{Phase 2: Fill remaining}
  \State $x^* \gets \argmax_{x \notin S} \min_{s \in S} \|x - s\|$
  \State $S \gets S \cup \{x^*\}$
\EndWhile
\State \Return $S$
\end{algorithmic}
\end{algorithm}

\section{Proof Details}
\label{app:submodular-proof}

\subsection{$(1-1/e)$ Approximation Bound}

\begin{theorem}[\citet{nemhauser1978analysis}]
Let $f: 2^V \to \mathbb{R}_{\geq 0}$ be a monotone submodular function.
The greedy algorithm that iteratively selects
$e^* = \argmax_{e \notin S} [f(S \cup \{e\}) - f(S)]$
achieves $f(S_k) \geq (1 - 1/e) \cdot f(S^*_k)$ where $S^*_k$
is the optimal size-$k$ subset.
\end{theorem}

\begin{proof}
Let $S^* = \{o_1, \ldots, o_k\}$ be optimal and $S_t$ the greedy
solution after $t$ steps. By monotonicity and submodularity:
\begin{align}
f(S^*) &\leq f(S_t \cup S^*) \leq f(S_t) + \sum_{j=1}^k [f(S_t \cup \{o_j\}) - f(S_t)]
\end{align}
Since greedy picks the element with maximum marginal gain:
\begin{align}
f(S_{t+1}) - f(S_t) &\geq \frac{1}{k}[f(S^*) - f(S_t)]
\end{align}
Let $\delta_t = f(S^*) - f(S_t)$. Then $\delta_{t+1} \leq (1 - 1/k)\delta_t$,
so $\delta_k \leq (1-1/k)^k \delta_0 \leq e^{-1} f(S^*)$, giving
$f(S_k) \geq (1 - 1/e)f(S^*)$.
\end{proof}

\subsection{Sum-Pairwise-Distance Properties}

The sum-pairwise-distance function $f(S) = \sum_{i<j,\, i,j\in S} d(i,j)$
is always monotone for non-negative distances. It is \emph{not}
submodular in general (consider three collinear points where adding
the middle point to a pair increases total distance more than adding
it to a singleton). However, in high-dimensional spaces with
near-uniform distributions, the function empirically behaves like
a submodular function, and the greedy algorithm performs well.

Our experimental results (Table~\ref{tab:baselines}) show the greedy
algorithm consistently outperforms random selection and K-Means
by large margins, suggesting that the $(1-1/e)$ bound provides a
useful, if conservative, quality guarantee in practice.

\section{Extended Experimental Results}
\label{app:extended}

\subsection{Full Selection Benchmark Across $k$ Values}

\begin{table}[H]
\centering
\small
\caption{Full selection benchmark on real text embeddings
(200 texts, $d=1536$) for $k \in \{5, 10, 20\}$.}
\label{tab:full_benchmark}
\begin{tabular}{llrrr}
\toprule
$k$ & Method & Spread & Sum Dist \\
\midrule
5 & FarthestPoint & $1.156 \pm 0.015$ & $12.4$ \\
5 & Submodular    & $1.128 \pm 0.000$ & $12.5$ \\
5 & MMR           & $1.066 \pm 0.000$ & $11.2$ \\
5 & DPP           & $0.656 \pm 0.226$ & $10.4$ \\
5 & Random        & $0.425 \pm 0.190$ & $9.0$ \\
\midrule
10 & FarthestPoint & $1.051 \pm 0.024$ & $53.7$ \\
10 & Submodular    & $0.902 \pm 0.000$ & $54.8$ \\
10 & MMR           & $0.713 \pm 0.000$ & $47.5$ \\
10 & DPP           & $0.442 \pm 0.122$ & $46.8$ \\
10 & Random        & $0.205 \pm 0.160$ & $39.5$ \\
\midrule
20 & FarthestPoint & $0.897 \pm 0.010$ & $221.0$ \\
20 & Submodular    & $0.455 \pm 0.000$ & $223.4$ \\
20 & MMR           & $0.499 \pm 0.000$ & $191.1$ \\
20 & DPP           & $0.361 \pm 0.063$ & $196.4$ \\
20 & Random        & $0.052 \pm 0.105$ & $163.4$ \\
\bottomrule
\end{tabular}
\end{table}

\paragraph{Analysis.}
As $k$ increases, the advantage of principled selection grows:
at $k=20$, FarthestPoint achieves spread $0.897$ vs.\ Random $0.052$
(17$\times$ improvement), and Submodular achieves sum distance $223.4$
vs.\ Random $163.4$ (37\% improvement). This demonstrates that
the toolkit is most valuable precisely when the selection problem
is hardest.

\subsection{Metric Taxonomy: Full GPT-2 Correlation Matrix}

The full $11 \times 11$ Kendall $\tau$ matrix on GPT-2 data:

{\small
\begin{table}[H]
\centering
\caption{GPT-2 Kendall $\tau$ matrix (15 temperature settings).}
\begin{tabular}{l|rrrrr}
\toprule
 & D-2 & SB & Ent & EPD & VS \\
\midrule
D-2  & 1.00 & $-0.90$ & 0.77 & 0.90 & 0.92 \\
SB   & $-0.90$ & 1.00 & $-0.75$ & $-0.92$ & $-0.98$ \\
Ent  & 0.77 & $-0.75$ & 1.00 & 0.68 & 0.73 \\
EPD  & 0.90 & $-0.92$ & 0.68 & 1.00 & 0.94 \\
VS   & 0.92 & $-0.98$ & 0.73 & 0.94 & 1.00 \\
\bottomrule
\end{tabular}

\medskip

\begin{tabular}{l|rrrrrr}
\toprule
 & Jac & PTD & CRD & USR & TTR & SED \\
\midrule
D-2  & 0.83 & 0.14 & 0.92 & 0.94 & 0.94 & 0.50 \\
SB   & $-0.89$ & $-0.09$ & $-0.90$ & $-0.96$ & $-0.92$ & $-0.45$ \\
PTD  & 0.16 & 1.00 & 0.18 & 0.12 & 0.14 & 0.22 \\
EPD  & 0.81 & 0.12 & 0.90 & 0.92 & 0.89 & 0.49 \\
\bottomrule
\end{tabular}
\end{table}
}

PTD is the clear outlier on GPT-2 ($|\tau| < 0.3$ with all other metrics),
confirming its independence. On \texttt{gpt-4.1-nano}, PTD collapses
into the lexical cluster, suggesting that smaller models produce
less syntactically diverse outputs.

\subsection{EPD Sample-Size Ablation}

\begin{table}[H]
\centering
\caption{EPD correlation with D-2 as a function of sample size
per group (\texttt{gpt-4.1-nano}).}
\begin{tabular}{rrr}
\toprule
Samples per group & Mean $\tau$(D-2, EPD) \\
\midrule
5  & $0.330$ \\
8  & $0.467$ \\
10 & $0.531$ \\
\bottomrule
\end{tabular}
\end{table}

EPD requires at least 8 samples per group to achieve stable
correlation estimates.

\subsection{Length Ablation}

\begin{table}[H]
\centering
\caption{EPD correlation with D-2 across generation lengths
(\texttt{gpt-4.1-nano}).}
\begin{tabular}{rrr}
\toprule
Max tokens & Mean $\tau$(D-2, EPD) \\
\midrule
30  & $0.503$ \\
50  & $0.533$ \\
100 & $0.531$ \\
\bottomrule
\end{tabular}
\end{table}

The EPD--D-2 correlation is stable across generation lengths,
indicating that the metric taxonomy is not an artifact of text length.

\subsection{Scaling Analysis Details}

\begin{table}[H]
\centering
\caption{Wall-clock time (ms) for $k=10$ selection across dataset sizes.}
\begin{tabular}{rrrrr}
\toprule
$n$ & DPP & Submod. & FarPt & Random \\
\midrule
50   & 1.4 & 3.0 & 0.4 & 0.2 \\
100  & 2.1 & 6.2 & 0.5 & 0.2 \\
200  & 6.3 & 16.0 & 0.6 & 0.2 \\
500  & 38.7 & 74.9 & 0.9 & 0.2 \\
1000 & 255.7 & 257.9 & 1.3 & 0.2 \\
2000 & -- & 1888.0 & 2.1 & 0.2 \\
\bottomrule
\end{tabular}
\end{table}

DPP becomes impractical above $n = 1000$ due to the $O(n^3)$
eigendecomposition. Submodular greedy scales as $O(n^{1.7})$
and remains usable up to $n = 2000$. FarthestPoint is the
fastest principled method at all scales.

\subsection{Fair Diversity Extended Analysis}

The fair selection algorithm (Section~\ref{sec:fair}) was tested
with imbalanced groups:
\begin{itemize}
\item Group 0 (60\%): 120 points, concentrated around a single centroid
\item Group 1 (25\%): 50 points, spread across a wider region
\item Group 2 (10\%): 20 points
\item Group 3 (5\%): 10 points
\end{itemize}

Without fairness constraints, FarthestPoint selects 2 from Group~0,
11 from Group~1, 4 from Group~2, and 3 from Group~3---overrepresenting
Group~1 (which has the most spread-out distribution). With
$\text{min}_g = 2$, the algorithm ensures at least 2 from each group,
including the concentrated Group~0, while maintaining 99.2\% of
total pairwise distance.

\section{Implementation Details}
\label{app:implementation}

\subsection{Unified Selector API}

All selectors implement the \texttt{DiversitySelector} abstract class:

\begin{verbatim}
class DiversitySelector(ABC):
    @abstractmethod
    def select(self, candidates: np.ndarray,
               k: int, **kwargs
    ) -> Tuple[List[int], Dict[str, Any]]:
        ...
\end{verbatim}

The factory function \texttt{get\_selector(name, **kwargs)} instantiates
any selector by name:

\begin{verbatim}
from divflow import get_selector
# All equivalent:
sel = get_selector('farthest_point')
sel = get_selector('dpp', kernel='rbf')
sel = get_selector('submodular', method='greedy')
sel = get_selector('mmr', lambda_param=0.3)
sel = get_selector('clustering')
sel = get_selector('random')
\end{verbatim}

\subsection{Numerical Stability}

DPP sampling requires eigendecomposition of the kernel matrix $L$.
We clamp eigenvalues to $[\epsilon, \infty)$ with $\epsilon = 10^{-10}$
to avoid numerical issues with near-singular matrices. For the
submodular optimizer, we precompute the full distance matrix to
avoid redundant distance computations during greedy selection.

\subsection{Dependencies}

\textsc{DivFlow} requires only \texttt{numpy} and \texttt{scipy}
(and optionally \texttt{scikit-learn} for comparison experiments).
This minimal dependency footprint facilitates adoption.

\section{Metric Definitions}
\label{app:metrics}

For completeness, we define all 11 diversity metrics.
Let $\mathcal{T} = \{t_1, \ldots, t_n\}$ be a set of texts.

\begin{itemize}
\item \textbf{Distinct-$n$} (D-$n$): $\frac{|\text{unique $n$-grams in } \mathcal{T}|}{|\text{total $n$-grams in } \mathcal{T}|}$

\item \textbf{Self-BLEU} (SB): Mean BLEU of each text against others:
$\text{SB} = \frac{1}{n} \sum_i \text{BLEU}(t_i, \mathcal{T} \setminus \{t_i\})$

\item \textbf{Type-Token Ratio} (TTR): $\frac{|\text{unique tokens}|}{|\text{total tokens}|}$

\item \textbf{Unique Sentence Ratio} (USR): Fraction of unique sentences
(measured by character 4-gram Jaccard similarity $< 0.5$).

\item \textbf{Compression Ratio Diversity} (CRD): $1 - \frac{|\text{compress}(\text{concat}(\mathcal{T}))|}{|\text{concat}(\mathcal{T})|}$.
Higher values indicate more diverse (less compressible) text.

\item \textbf{$n$-gram Entropy} (Ent): $H = -\sum_g p(g) \log_2 p(g)$
over all $n$-grams $g$ in $\mathcal{T}$.

\item \textbf{Jaccard Distance} (Jac): $1 - \frac{1}{\binom{n}{2}} \sum_{i<j} J(t_i, t_j)$
where $J$ is bigram Jaccard similarity.

\item \textbf{Parse-Tree Diversity} (PTD): Mean pairwise tree-edit distance
of syntactic parse trees.

\item \textbf{Semantic Embedding Diversity} (SED): Mean pairwise cosine
distance in sentence embedding space.

\item \textbf{Embedding Pairwise Distance} (EPD): Mean pairwise Euclidean
distance in embedding space.

\item \textbf{Vendi Score} (VS): $\exp(H(\lambda_1, \ldots, \lambda_n))$
where $\lambda_i$ are eigenvalues of the normalised similarity matrix
and $H$ is Shannon entropy.
\end{itemize}

\section{Artifact Traceability}
\label{app:traceability}

Every quantitative claim in this paper maps to a specific JSON artifact:

\begin{table}[H]
\centering
\small
\caption{Claim-to-artifact mapping.}
\begin{tabular}{p{6cm}p{5.5cm}}
\toprule
Claim & Artifact \\
\midrule
D-2 vs SB $\tau = -0.947$ (nano) & \texttt{unified\_results.json} $\to$ \texttt{nano.kendall\_tau\_per\_prompt} \\
10/11 metrics cluster at $|\tau|>0.7$ & \texttt{unified\_results.json} $\to$ \texttt{nano.clusters\_at\_0.7} \\
EPD $\tau = 0.531$ with D-2 & \texttt{unified\_results.json} $\to$ \texttt{nano.kendall\_tau\_per\_prompt} \\
FarthestPoint spread 10.64 & \texttt{comprehensive\_results.json} $\to$ \texttt{baseline\_comparison} \\
Submodular sum\_dist 547.5 & \texttt{comprehensive\_results.json} $\to$ \texttt{baseline\_comparison} \\
FP vs Random $d = 5.93$ & \texttt{comprehensive\_results.json} $\to$ \texttt{stat\_test\_fp\_vs\_random} \\
Fair diversity retention 99.2\% & \texttt{comprehensive\_results.json} $\to$ \texttt{fair\_diversity} \\
D-2 improvement 52\% & \texttt{comprehensive\_results.json} $\to$ \texttt{text\_selection} \\
Scaling exponents & \texttt{comprehensive\_results.json} $\to$ \texttt{scaling} \\
\bottomrule
\end{tabular}
\end{table}

\end{document}
